\subsection*{سؤال 15 امتحان 899205 سنة 2011}
\addcontentsline{toc}{subsection}{سؤال 15 امتحان 899205 سنة 2011}

\insertFullImg{../../../bagrut_questions/computational_models/irregularity_2011_899205_15.png}

\noindent \large{ج. ابن أوتوماتًا نهائيَّا محدودًا يقبل اللغة $L_1$.}

\ifwithsols
\begin{boxSolution}[سؤال 15 امتحان 899205 سنة 2011]
\begin{enumerate}[itemsep=1.5em, label=\alph*.]
\item نبرهن ذلك بالتناقض: نفترض العكس: نفترض أن $L$ نظامية. ونثبت أن هذا يؤدي إلى تناقض. \\
إذا كانت $L$ نظامية فهذا يعني أنّ هناك أوتوماتًا نهائيًّا محدودًا $A$ يقبلها. نفترض أنّ $A$ مكون من $m$ حالات مختلفة. \\
نعرّف مجموعة الكلمات التالية: $B = \{0^p \mid p > 0 \}$، من الواضح أنّ $B$ هي مجموعة لا نهائية. لذلك، وحسب مبدأ برج الحمام، بالضرورة هناك عددان صحيحان $i > j$ بحيث إذا شغّلنا الأوتومات $A$ على الكلمتين $w_i=0^i$ و $w_j=0^j$ فإنّهما ستتوقفان في نفس الحالة $q$. \\
ننظر إلى الكلمتين: $w_i' = 0^i 1^j 2$ و $w_j' = 0^j 1^j 2$. \\
من الواضح حسب تعريف اللغة أنّ $w_i' \in L$ بينما $w_j' \notin L$.
لكن بما أنّ تشغيل الأوتومات على المقطعين $0^i$ و $0^j$ ينتهي بنفس الحالة $q$ فإنّ الاستمرار من هناك بنفس اللاحقة $1^j 2$ يعني أنّ الكلمتين $w_i'$ و $w_j'$ ستنتهيان أيضًا في نفس الحالة، أي إمّا أنّ كليهما مقبولتان أو كلاهما مرفوضتان، وهذا يناقض أنّ $A$ يقبل اللغة $L$.

\item $L_1 \cap L2 = \emptyset$ \\
لأنّ كل الكلمات في $L$ تبدأ بالحرف $0$، بينما في $L_1$ فإنّه لا توجد كلمات تحتوي على $0$ أصلًا.

\item .
\begin{center}
    \includegraphics[width=0.4\textwidth]{../../../images/DFAs/1power_n_2.png}
\end{center}
\end{enumerate}
\end{boxSolution}
\clearpage
\fi
